% GENERATED FILE. Edit canonical lessons, then run npm run generate:books.
% canonical-source-hash: fnv1a64:55affd60edcf3666
% canonical-lessons: MR-C451-prasanna, MR-C451-udas, MR-C451-khus, MR-C451-naraj, MR-C451-vaitagne

\chapter{Cheerful and Gloomy}
\label{ch:mr-cheerful-and-gloomy}
\hlchaptermodality{451}

\begin{chapteropening}
\textbf{By the end of this chapter:} \emph{I can say whether someone is cheerful, gloomy, pleased or upset.}

Complete the last lesson of chapter 451: I can say whether someone is cheerful, gloomy, pleased or upset.
\end{chapteropening}

\section[\texorpdfstring{\mr{प्रसन्न} (prasanna)}{प्रसन्न (prasanna)}]{\mr{प्रसन्न} (prasanna) — cheerful; bright and clear}

\label{lesson:MR-C451-prasanna}



\begin{quote}
Before the new one: say the Marathi for surprise, then the Marathi for to be deeply moved.
\end{quote}

\subsection*{You'll want to know: \mr{प्रसन्न}}
\textbf{\mr{प्रसन्न}} — \emph{prasanna} — \textquotedblleft{}cheerful; bright and clear\textquotedblright{}.

\begin{quote}
\mr{सकाळी} \mr{बागेत} \mr{फिरल्यावर} \mr{मन} \mr{प्रसन्न} \mr{होतं}.
\end{quote}

\begin{quote}
\emph{sakāḷī bāget phirlyāvar man prasanna hotaṁ.}
\end{quote}

\begin{quote}
After a morning walk in the garden, the mind feels bright.
\end{quote}

Unchanging for gender: \mr{प्रसन्न} \mr{चेहरा}, a cheerful face; \mr{प्रसन्न} \mr{सकाळ}, a bright morning.

\begin{cousinweb}
\textbf{\mr{प्रसन्न}} (\emph{prasanna}) is a tatsama, taken from Sanskrit \emph{prasanna}, \textquotedblleft{}pleased; clear, bright\textquotedblright{} (CDIAL 8823).
\end{cousinweb}

\begin{grammarlens}[title={What you've built}]
One more word for describing what you lived and felt.
\end{grammarlens}

\subsection*{Your turn}
\begin{itemize}
\raggedright
  \item \emph{Say it:} \emph{prasanna}
  \item \emph{Say it:} \emph{prasanna}, once more
  \item \emph{Say it:} \emph{sakāḷī bāget phirlyāvar man prasanna hotaṁ.} — After a morning walk in the garden, the mind feels bright.
  \item \emph{Recall:} say the Marathi for surprise, then the Marathi for to be deeply moved, then say \emph{prasanna} again
\end{itemize}

\begin{tcolorbox}[breakable,colback=teal!4,colframe=teal!35!black,title={Before you move on}]
What does \mr{प्रसन्न} mean? (\textquotedblleft{}Cheerful; bright and clear\textquotedblright{}.) Say it once more.
\end{tcolorbox}
\section[\texorpdfstring{\mr{उदास} (udās)}{उदास (udās)}]{\mr{उदास} (udās) — gloomy; low, downcast}

\label{lesson:MR-C451-udas}



\begin{quote}
Before the new one: say the Marathi for to be deeply moved, then the Marathi for cheerful.
\end{quote}

\subsection*{You'll want to know: \mr{उदास}}
\textbf{\mr{उदास}} — \emph{udās} — \textquotedblleft{}gloomy; low, downcast\textquotedblright{}.

\begin{quote}
\mr{पावसाळ्यात} \mr{अरुण} \mr{कधीकधी} \mr{उदास} \mr{असतो}.
\end{quote}

\begin{quote}
\emph{pāvsāḷyāt aruṇ kadhīkadhī udās asto.}
\end{quote}

\begin{quote}
In the monsoon Arun is sometimes low.
\end{quote}

Unchanging: \mr{उदास} \mr{मुलगा}, \mr{उदास} \mr{मुलगी}. \mr{उदास} \mr{वाटणे}, to feel low.

\begin{cousinweb}
\textbf{\mr{उदास}} (\emph{udās}) is usually traced to Sanskrit \emph{udāsa}, a noun, \textquotedblleft{}indifference, apathy\textquotedblright{} (Monier-Williams); in Marathi it is an adjective of low spirits.
\end{cousinweb}

\begin{grammarlens}[title={What you've built}]
One more word for describing what you lived and felt.
\end{grammarlens}

\subsection*{Your turn}
\begin{itemize}
\raggedright
  \item \emph{Say it:} \emph{udās}
  \item \emph{Say it:} \emph{udās}, once more
  \item \emph{Say it:} \emph{pāvsāḷyāt aruṇ kadhīkadhī udās asto.} — In the monsoon Arun is sometimes low.
  \item \emph{Recall:} say the Marathi for to be deeply moved, then the Marathi for cheerful, then say \emph{udās} again
\end{itemize}

\begin{tcolorbox}[breakable,colback=teal!4,colframe=teal!35!black,title={Before you move on}]
What does \mr{उदास} mean? (\textquotedblleft{}Gloomy; low, downcast\textquotedblright{}.) Say it once more.
\end{tcolorbox}
\section[\texorpdfstring{\mr{खुश} (khuś)}{खुश (khuś)}]{\mr{खुश} (khuś) — pleased; glad}

\label{lesson:MR-C451-khus}



\begin{quote}
Before the new one: say the Marathi for cheerful, then the Marathi for gloomy.
\end{quote}

\subsection*{You'll want to know: \mr{खुश}}
\textbf{\mr{खुश}} — \emph{khuś} — \textquotedblleft{}pleased; glad\textquotedblright{}.

\begin{quote}
\mr{बक्षीस} \mr{मिळाल्यावर} \mr{मुली} \mr{खुश} \mr{झाल्या}.
\end{quote}

\begin{quote}
\emph{bakṣīs miḷālyāvar mulī khuś jhālyā.}
\end{quote}

\begin{quote}
The girls were pleased when they got the prize.
\end{quote}

Unchanging; \mr{खुश} \mr{होणे}, to become pleased. A Persian word beside the Sanskrit \mr{आनंदी}.

\begin{cousinweb}
\textbf{\mr{खुश}} (\emph{khuś}) is a loan from Persian \emph{khush}, \textquotedblleft{}happy, pleasant\textquotedblright{}; Marathi writes the Persian \emph{kh} with \mr{ख}.
\end{cousinweb}

\begin{grammarlens}[title={What you've built}]
One more word for describing what you lived and felt.
\end{grammarlens}

\subsection*{Your turn}
\begin{itemize}
\raggedright
  \item \emph{Say it:} \emph{khuś}
  \item \emph{Say it:} \emph{khuś}, once more
  \item \emph{Say it:} \emph{bakṣīs miḷālyāvar mulī khuś jhālyā.} — The girls were pleased when they got the prize.
  \item \emph{Recall:} say the Marathi for cheerful, then the Marathi for gloomy, then say \emph{khuś} again
\end{itemize}

\begin{tcolorbox}[breakable,colback=teal!4,colframe=teal!35!black,title={Before you move on}]
What does \mr{खुश} mean? (\textquotedblleft{}Pleased; glad\textquotedblright{}.) Say it once more.
\end{tcolorbox}
\section[\texorpdfstring{\mr{नाराज} (nārāj)}{नाराज (nārāj)}]{\mr{नाराज} (nārāj) — displeased; upset}

\label{lesson:MR-C451-naraj}



\begin{quote}
Before the new one: say the Marathi for gloomy, then the Marathi for pleased.
\end{quote}

\subsection*{You'll want to know: \mr{नाराज}}
\textbf{\mr{नाराज}} — \emph{nārāj} — \textquotedblleft{}displeased; upset\textquotedblright{}.

\begin{quote}
\mr{मी} \mr{उशिरा} \mr{आलो} \mr{म्हणून} \mr{वडील} \mr{नाराज} \mr{झाले}.
\end{quote}

\begin{quote}
\emph{mī uśirā ālo mhaṇūn vaḍīl nārāj jhāle.}
\end{quote}

\begin{quote}
Father was displeased because I came late.
\end{quote}

Unchanging. The person one is upset with takes -\mr{वर}: \mr{ती} \mr{माझ्यावर} \mr{नाराज} \mr{आहे}.

\begin{cousinweb}
\textbf{\mr{नाराज}} (\emph{nārāj}) is a loan from Persian \emph{nārāz}, \textquotedblleft{}displeased\textquotedblright{}, built with the Persian negative \emph{nā-}.
\end{cousinweb}

\begin{grammarlens}[title={What you've built}]
One more word for describing what you lived and felt.
\end{grammarlens}

\subsection*{Your turn}
\begin{itemize}
\raggedright
  \item \emph{Say it:} \emph{nārāj}
  \item \emph{Say it:} \emph{nārāj}, once more
  \item \emph{Say it:} \emph{mī uśirā ālo mhaṇūn vaḍīl nārāj jhāle.} — Father was displeased because I came late.
  \item \emph{Recall:} say the Marathi for gloomy, then the Marathi for pleased, then say \emph{nārāj} again
\end{itemize}

\begin{tcolorbox}[breakable,colback=teal!4,colframe=teal!35!black,title={Before you move on}]
What does \mr{नाराज} mean? (\textquotedblleft{}Displeased; upset\textquotedblright{}.) Say it once more.
\end{tcolorbox}
\section[\texorpdfstring{\mr{वैतागणे} (vaitāgṇe)}{वैतागणे (vaitāgṇe)}]{\mr{वैतागणे} (vaitāgṇe) — to get fed up; to be exasperated}

\label{lesson:MR-C451-vaitagne}



\begin{quote}
Before the new one: say the Marathi for pleased, then the Marathi for displeased.
\end{quote}

\subsection*{You'll want to know: \mr{वैतागणे}}
\textbf{\mr{वैतागणे}} — \emph{vaitāgṇe} — \textquotedblleft{}to get fed up; to be exasperated\textquotedblright{}.

\begin{quote}
\mr{दिवसभर} \mr{पाऊस} \mr{पडत} \mr{होता}, \mr{म्हणून} \mr{मुलं} \mr{वैतागली}.
\end{quote}

\begin{quote}
\emph{divasbhar pāūs paḍat hotā, mhaṇūn mulaṁ vaitāglī.}
\end{quote}

\begin{quote}
It rained all day, so the children got fed up.
\end{quote}

Past (he …): \mr{तो} \mr{वैतागला}, \emph{to vaitāglā}.

What one is fed up with takes -\mr{ला}: \mr{कामाला} \mr{वैतागणे}. The noun \mr{वैताग}, exasperation.

\begin{cousinweb}
\textbf{\mr{वैतागणे}} (\emph{vaitāgṇe}) is a Marathi verb built on \mr{वैताग}, \textquotedblleft{}exasperation\textquotedblright{}; the origin of \mr{वैताग} is not settled.
\end{cousinweb}

\begin{grammarlens}[title={What you've built}]
One more word for describing what you lived and felt.
\end{grammarlens}

\subsection*{Your turn}
\begin{itemize}
\raggedright
  \item \emph{Say it:} \emph{vaitāgṇe}
  \item \emph{Say it:} \emph{vaitāgṇe}, once more
  \item \emph{Say it:} \emph{divasbhar pāūs paḍat hotā, mhaṇūn mulaṁ vaitāglī.} — It rained all day, so the children got fed up.
  \item \emph{Recall:} say the Marathi for pleased, then the Marathi for displeased, then say \emph{vaitāgṇe} again
\end{itemize}

\begin{tcolorbox}[breakable,colback=teal!4,colframe=teal!35!black,title={Before you move on}]
What does \mr{वैतागणे} mean? (\textquotedblleft{}To get fed up; to be exasperated\textquotedblright{}.) Say it once more.
\end{tcolorbox}
